% Paper 2 	extperiodcentered{}  §2 Related work
% Status: draft 	extperiodcentered{}  2026-05-05

\section{Related work}
\label{sec:related}

\subsection{Long-term memory benchmarks}

LongMemEval-S \citep{wu2024longmemeval} is the standard public benchmark
for LLM long-term memory recall, defining 500 questions across 6 cognitive
abilities. Prior to LongMemEval-S, memory was evaluated indirectly through
multi-turn dialog quality \citep{ouyang2022training} or via proxy retrieval
metrics like recall@k \citep{karpukhin2020dpr}. The released subset
(LongMemEval-S, 500 questions, 50K-token haystacks) standardizes the
evaluation and permits direct comparison.

\subsection{Memory systems}

\textbf{Mem0} \citep{mem0_2024} retrieves memory using LLM-summarized
embeddings; benchmarks at 40--45\% on LongMemEval-S. Heavy LLM
involvement at retrieval time is expensive and slow.

\textbf{Letta} (formerly MemGPT) \citep{letta2024} hierarchical memory
with manual archival; 35--38\% on LongMemEval-S. Truncates rather than
prioritizes when context is full.

\textbf{A-MEM} \citep{amem_2024} adaptive memory with learned forgetting;
\(\sim\)50\%. Trains a forget-or-keep classifier per memory.

\textbf{Zep} \citep{zep2025} knowledge-graph memory using entity extraction
and relationship traversal; 55--60\% (top of SOTA). Requires Neo4j and
ongoing graph maintenance.

\textbf{claude-mem} \citep{claudemem2026} narrative-summary plugin for
Claude Code; LLM rewrites session into structured observation. Focuses on
narrative quality rather than retrieval accuracy. Does not benchmark on
LongMemEval-S.

In contrast, Compass v0.8 is the first to combine: (1) a public
500-question benchmark result, (2) sub-1/15 cost compared to Zep, and (3)
multi-protocol integration (MCP + A2A) without requiring graph
infrastructure.

\subsection{Persona drift and self-audit}

Anthropic's \emph{Persona Vectors} paper \citep{anthropic2025persona}
demonstrated that LLM activations contain steerable directions for
sycophancy and hallucination, but the technique requires white-box
access. Black-box equivalents using prompt-based detection
\citep{wang2024selfdetect} achieve weaker AUC (\(<\)0.7) and high false
positive rates.

Compass's anchor-based drift detection (\S\ref{sec:method}.A appendix)
uses 25 positive + 35 negative anchor sentences, embedded with bge-m3 and
compared via cosine to incoming prompts. We report AUC=0.92 on a
held-out 200-prompt test set—the first black-box drift detector
suitable for production hooks at \(<\)50\,ms p95 latency. This contribution
is orthogonal to the LongMemEval-S pipeline but informs the design of
the MCP server's \texttt{drift\_history} tool.

\subsection{LLM thinking-mode behaviors}

OpenAI's o1 \citep{openai2024o1}, DeepSeek's r1 \citep{deepseek2025r1}, and
Anthropic's extended thinking \citep{anthropic2024thinking} introduced
explicit chain-of-thought tokens. Recent work \citep{wang2025thinkingharm}
showed thinking can hurt accuracy on tasks where the LLM's prior is
already correct (over-reasoning effect). Our ablation extends this:
thinking helps DeepSeek V3.2 on temporal-reasoning (\(+6.8\) pts) but
shows zero gain for Kimi K2.6 and triggers refusal cascade in MiniMax
M2.7-highspeed at low budget. The MiniMax cascade is the strongest
documented case of thinking causing systematic failure that we are
aware of, and warrants inclusion in any production deployment
checklist.

\subsection{Cross-agent memory federation}

To our knowledge, no prior public system addresses
\emph{cross-agent} memory sharing—where the same user's interactions
across Claude Desktop, Cursor, Cline, OpenClaw, etc., contribute to a
unified memory. Cross-platform clipboard and password managers (1Password,
iCloud Keychain) provide analogous single-user-multi-device unification
but for static data, not interaction logs. Compass v0.9 introduces
the multi-agent ingest API (\texttt{POST /v1/observations} keyed on
\texttt{user\_id} + \texttt{agent\_id}) and ships SDK adapters for the
five major MCP-compatible clients. We are not aware of competing
implementations.
